\documentclass[10pt,a4paper]{amsart}
\usepackage{amsmath,amssymb,amsthm,graphicx}
% For commutative diagrams we use Xy-pic
\usepackage[all]{xy}

% We use 2cell for 2-commutative diagrams.
\xyoption{2cell}
% Theorem environments.
%
\theoremstyle{plain}
\newtheorem{theorem}[subsection]{Theorem}
\newtheorem{proposition}[subsection]{Proposition}
\newtheorem{lemma}[subsection]{Lemma}

\theoremstyle{definition}
\newtheorem{definition}[subsection]{Definition}
\newtheorem{example}[subsection]{Example}
\newtheorem{exercise}[subsection]{Exercise}
\newtheorem{situation}[subsection]{Situation}

\theoremstyle{remark}
\newtheorem{remark}[subsection]{Remark}
\newtheorem{remarks}[subsection]{Remarks}

\numberwithin{equation}{subsection}

% Macros
%
\def\lim{\mathop{\mathrm{lim}}\nolimits}
\def\colim{\mathop{\mathrm{colim}}\nolimits}
\def\Spec{\mathop{\mathrm{Spec}}}
\def\Hom{\mathop{\mathrm{Hom}}\nolimits}
\def\Ext{\mathop{\mathrm{Ext}}\nolimits}
\def\SheafHom{\mathop{\mathcal{H}\!\mathit{om}}\nolimits}
\def\SheafExt{\mathop{\mathcal{E}\!\mathit{xt}}\nolimits}
\def\Sch{\mathit{Sch}}
\def\Mor{\mathop{\mathrm{Mor}}\nolimits}
\def\Ob{\mathop{\mathrm{Ob}}\nolimits}
\def\Sh{\mathop{\mathit{Sh}}\nolimits}
\def\NL{\mathop{N\!L}\nolimits}
\def\CH{\mathop{\mathrm{CH}}\nolimits}
\def\proetale{{pro\text{-}\acute{e}tale}}
\def\etale{{\acute{e}tale}}
\def\QCoh{\mathit{QCoh}}
\def\Ker{\mathop{\mathrm{Ker}}}
\def\Im{\mathop{\mathrm{Im}}}
\def\Coker{\mathop{\mathrm{Coker}}}
\def\Coim{\mathop{\mathrm{Coim}}}

% Boxtimes
%
\DeclareMathSymbol{\boxtimes}{\mathbin}{AMSa}{"02}

%
% Macros for moduli stacks/spaces
%
\def\QCohstack{\mathcal{QC}\!\mathit{oh}}
\def\Cohstack{\mathcal{C}\!\mathit{oh}}
\def\Spacesstack{\mathcal{S}\!\mathit{paces}}
\def\Quotfunctor{\mathrm{Quot}}
\def\Hilbfunctor{\mathrm{Hilb}}
\def\Curvesstack{\mathcal{C}\!\mathit{urves}}
\def\Polarizedstack{\mathcal{P}\!\mathit{olarized}}
\def\Complexesstack{\mathcal{C}\!\mathit{omplexes}}
% \Pic is the operator that assigns to X its picard group, usage \Pic(X)
% \Picardstack_{X/B} denotes the Picard stack of X over B
% \Picardfunctor_{X/B} denotes the Picard functor of X over B
\def\Pic{\mathop{\mathrm{Pic}}\nolimits}
\def\Picardstack{\mathcal{P}\!\mathit{ic}}
\def\Picardfunctor{\mathrm{Pic}}
\def\Deformationcategory{\mathcal{D}\!\mathit{ef}}
\usepackage{cleveref}
\crefname{originallemma}{original lemma}{original lemmas}
\Crefname{originallemma}{Original lemma}{Original lemmas}
\makeatletter
\newlabel{lemma-colimit-variant-complexes}{{0EZM}{source}{}{sourcepointer-1}{}}
\newlabel{lemma-colimit-variant-complexes@cref}{{[originallemma][0][]0EZM}{[1][0][]source}{}{}{}}
\makeatother
\title{Etale cohomology of a directed limit of quasi-compact quasi-separated schemes is the colimit of the cohomologies}
\author{The Stacks Project Authors}
\date{stacks-project:etale-cohomology.tex:local-snapshot-20261011; GNU FDL 1.2 or later}
\begin{document}
\maketitle
\noindent\textit{Editorial note.} All delivered mathematical text is the author's own, copied from the pinned Stacks Project chapter file etale-cohomology.tex; no statement or proof step was rewritten, reordered or removed.\par\medskip
\noindent\textit{Editorial note.} The proof is the author's complete argument: it supplies the representative complex, the pullback bookkeeping and the termwise colimit identity, and then cites the earlier variant lemma for the final identification.\par\medskip
\noindent\textit{Editorial note.} The single labelled reference occurring inside the counted spans is carried in assembly.external\_references as a non-load-bearing navigation pointer to the same chapter, so no load-bearing mathematics is replaced by a stub.\par\medskip
\noindent\textit{Editorial note.} The label lemma-colimit-complexes occurs in no other staged Stacks chapter, so no cross-chapter restatement had to be weighed for this card.\par\medskip
\noindent\textit{Editorial note.} The preamble is the chapter's own native preamble excerpt already used for this source: the Xy-pic setup (native preamble lines 14-18) and the native theorem declarations with the mathematical macros (native preamble lines 154-217).\par\medskip
\noindent\textit{Allowed standard prerequisites (editorial).}\begin{itemize}
\item Derived category D\textasciicircum{}+(S\_etale) and the pullback of a complex to a scheme over S. Same chapter, the bounded below derived category of abelian sheaves on S\_etale and the pullbacks K|\_\{X\_i\} and K|\_X. Supplies the object K whose cohomology is being compared and the notation for its pullbacks.
\item Representation of an object of D\textasciicircum{}+(S\_etale) by a bounded below complex of abelian sheaves. Same chapter, the description of D\textasciicircum{}+(S\_etale) by bounded below complexes. Supplies the complex G\textasciicircum{}bullet with G\textasciicircum{}n = 0 for n < a which the proof uses to replace K.
\item Identification of the pulled back complex as the termwise colimit of the pulled back complexes. Same chapter, the pullback of a complex along the transition morphisms and the termwise colimit description F\textasciicircum{}bullet = colim f\_i\textasciicircum{}\{-1\} F\_i\textasciicircum{}bullet. Supplies the identity which, together with the cited variant lemma, gives the statement.
\end{itemize}\medskip
\ifdefined\hypertarget\hypertarget{sourcepointer-1}{}\fi
\noindent\textit{Source reference.} Original lemma 0EZM: sources/stacks-etale-cohomology/pinned-original-source.tex, source lines 7783--7783. Same-chapter navigation pointer invoked at proof line 7873, where the author cites the earlier variant lemma on colimits of complexes. The pointer only resolves that reference for the reader; the cited lemma is not reproduced inside the counted span, and no bibliography citation is involved.\par\medskip
\begin{lemma}
\label{lemma-colimit-complexes}
Let $S$ be a scheme. Let $X = \lim_{i \in I} X_i$ be a limit of a
directed system of schemes over $S$ with affine transition morphisms
$f_{i'i} : X_{i'} \to X_i$. We assume that $X_i$ is quasi-compact and
quasi-separated for all $i \in I$. Let $K \in D^+(S_\etale)$. Then
$$
\colim_{i \in I} H_\etale^p(X_i, K|_{X_i}) = H_\etale^p(X, K|_X).
$$
for all $p \in \mathbf{Z}$ where $K|_{X_i}$ and $K|_X$
are the pullbacks of $K$ to $X_i$ and $X$.
\end{lemma}

\begin{proof}
We may represent $K$ by a bounded below complex $\mathcal{G}^\bullet$
of abelian sheaves on $S_\etale$. Say $\mathcal{G}^n = 0$ for $n < a$.
Denote $\mathcal{F}^\bullet_i$ and $\mathcal{F}^\bullet$ the pullbacks
of this complex of $X_i$ and $X$. These complexes represent the
objects $K|_{X_i}$ and $K|_X$ and we have
$\mathcal{F}^\bullet = \colim f_i^{-1}\mathcal{F}_i^\bullet$
termwise. Hence the lemma follows from
Lemma \ref{lemma-colimit-variant-complexes}.
\end{proof}

\end{document}
